\documentclass[12pt,a4paper]{exam}
\usepackage[utf8]{inputenc}
\usepackage{amsmath,amssymb,amsthm}
\usepackage[margin=1in]{geometry}
\usepackage{graphicx}
\usepackage[colorlinks=true]{hyperref}
\pagestyle{headandfoot}
\firstpageheader{MIT OpenCourseWare}{}{\textbf{EXTREMELY HARD -- MIT LEVEL}}
\runningheader{MIT OpenCourseWare}{}{\textbf{EXTREMELY HARD -- MIT LEVEL}}
\firstpagefooter{}{Page \thepage}{}
\runningfooter{}{Page \thepage}{}

\begin{document}

\begin{center}
{\LARGE \textbf{MIT OpenCourseWare}}\\21L.005\\Introduction to Drama
\vspace{0.5cm}
{21L.005} --- {Introduction to Drama}
\vspace{0.3cm}
May 2026
\vspace{0.3cm}
\textbf{EXTREMELY HARD -- MIT LEVEL}
\end{center}

\begin{center}
\textbf{Time Limit: 3 hours}
\end{center}

\begin{center}
\textbf{Total Marks: 100}
\end{center}

\vspace{0.5cm}

\begin{questions}

\question[3]
Prove the fundamental theorem concerning Greek Tragedy in the context of 21L.005 (Introduction to Drama). State the theorem precisely, provide a complete proof, and discuss three important corollaries.

\vspace{0.3cm}

\question[3]
Analyze the limitations of standard approaches to Shakespeare when applied to edge cases in 21L.005. Construct explicit counterexamples showing where intuitive reasoning fails.

\vspace{0.3cm}

\question[3]
Analyze the limitations of standard approaches to Modern Drama when applied to edge cases in 21L.005. Construct explicit counterexamples showing where intuitive reasoning fails.

\vspace{0.3cm}

\question[4]
Construct an original proof-level problem involving Performance Theory for 21L.005 (Introduction to Drama) and provide a complete solution. Your problem should be at the level of a graduate qualifying exam.

\vspace{0.3cm}

\question[6]
Analyze the limitations of standard approaches to Theatricality when applied to edge cases in 21L.005. Construct explicit counterexamples showing where intuitive reasoning fails.

\vspace{0.3cm}

\question[5]
Construct an original proof-level problem involving Greek Tragedy for 21L.005 (Introduction to Drama) and provide a complete solution. Your problem should be at the level of a graduate qualifying exam.

\vspace{0.3cm}

\question[3]
Analyze the limitations of standard approaches to Shakespeare when applied to edge cases in 21L.005. Construct explicit counterexamples showing where intuitive reasoning fails.

\vspace{0.3cm}

\question[4]
Construct an original proof-level problem involving Modern Drama for 21L.005 (Introduction to Drama) and provide a complete solution. Your problem should be at the level of a graduate qualifying exam.

\vspace{0.3cm}

\question[3]
Prove the fundamental theorem concerning Performance Theory in the context of 21L.005 (Introduction to Drama). State the theorem precisely, provide a complete proof, and discuss three important corollaries.

\vspace{0.3cm}

\question[3]
Construct an original proof-level problem involving Theatricality for 21L.005 (Introduction to Drama) and provide a complete solution. Your problem should be at the level of a graduate qualifying exam.

\vspace{0.3cm}

\question[3]
Construct an original proof-level problem involving Greek Tragedy for 21L.005 (Introduction to Drama) and provide a complete solution. Your problem should be at the level of a graduate qualifying exam.

\vspace{0.3cm}

\question[3]
Prove the fundamental theorem concerning Shakespeare in the context of 21L.005 (Introduction to Drama). State the theorem precisely, provide a complete proof, and discuss three important corollaries.

\vspace{0.3cm}

\question[7]
Prove the fundamental theorem concerning Modern Drama in the context of 21L.005 (Introduction to Drama). State the theorem precisely, provide a complete proof, and discuss three important corollaries.

\vspace{0.3cm}

\question[5]
Analyze the limitations of standard approaches to Performance Theory when applied to edge cases in 21L.005. Construct explicit counterexamples showing where intuitive reasoning fails.

\vspace{0.3cm}

\question[3]
Prove the fundamental theorem concerning Theatricality in the context of 21L.005 (Introduction to Drama). State the theorem precisely, provide a complete proof, and discuss three important corollaries.

\vspace{0.3cm}

\question[5]
Construct an original proof-level problem involving Greek Tragedy for 21L.005 (Introduction to Drama) and provide a complete solution. Your problem should be at the level of a graduate qualifying exam.

\vspace{0.3cm}

\question[3]
Construct an original proof-level problem involving Shakespeare for 21L.005 (Introduction to Drama) and provide a complete solution. Your problem should be at the level of a graduate qualifying exam.

\vspace{0.3cm}

\question[6]
Prove the fundamental theorem concerning Modern Drama in the context of 21L.005 (Introduction to Drama). State the theorem precisely, provide a complete proof, and discuss three important corollaries.

\vspace{0.3cm}

\question[6]
Prove the fundamental theorem concerning Performance Theory in the context of 21L.005 (Introduction to Drama). State the theorem precisely, provide a complete proof, and discuss three important corollaries.

\vspace{0.3cm}

\question[2]
Prove the fundamental theorem concerning Theatricality in the context of 21L.005 (Introduction to Drama). State the theorem precisely, provide a complete proof, and discuss three important corollaries.

\vspace{0.3cm}

\question[5]
Construct an original proof-level problem involving Greek Tragedy for 21L.005 (Introduction to Drama) and provide a complete solution. Your problem should be at the level of a graduate qualifying exam.

\vspace{0.3cm}

\question[3]
Analyze the limitations of standard approaches to Shakespeare when applied to edge cases in 21L.005. Construct explicit counterexamples showing where intuitive reasoning fails.

\vspace{0.3cm}

\question[5]
Prove the fundamental theorem concerning Modern Drama in the context of 21L.005 (Introduction to Drama). State the theorem precisely, provide a complete proof, and discuss three important corollaries.

\vspace{0.3cm}

\question[3]
Prove the fundamental theorem concerning Performance Theory in the context of 21L.005 (Introduction to Drama). State the theorem precisely, provide a complete proof, and discuss three important corollaries.

\vspace{0.3cm}

\question[4]
Analyze the limitations of standard approaches to Theatricality when applied to edge cases in 21L.005. Construct explicit counterexamples showing where intuitive reasoning fails.

\vspace{0.3cm}

\end{questions}

\newpage
\section*{Solutions}

\textbf{Solution 1:}

The proof follows from the definitions and properties established in 21L.005. The key insight is to apply the relevant theorem and verify the hypotheses hold. Detailed solution depends on the specific formulation of the theorem.

\vspace{0.5cm}

\textbf{Solution 2:}

The edge case analysis reveals the subtle assumptions underlying standard treatments of Shakespeare. By constructing explicit counterexamples, we see that the usual hypotheses cannot be weakened without loss of the theorem's conclusion.

\vspace{0.5cm}

\textbf{Solution 3:}

The edge case analysis reveals the subtle assumptions underlying standard treatments of Modern Drama. By constructing explicit counterexamples, we see that the usual hypotheses cannot be weakened without loss of the theorem's conclusion.

\vspace{0.5cm}

\textbf{Solution 4:}

Solution approach: First recall the definitions and key theorems related to Performance Theory from 21L.005. Then apply the appropriate techniques to construct a rigorous argument. The solution must be self-contained and reference only the material from the course.

\vspace{0.5cm}

\textbf{Solution 5:}

The edge case analysis reveals the subtle assumptions underlying standard treatments of Theatricality. By constructing explicit counterexamples, we see that the usual hypotheses cannot be weakened without loss of the theorem's conclusion.

\vspace{0.5cm}

\textbf{Solution 6:}

Solution approach: First recall the definitions and key theorems related to Greek Tragedy from 21L.005. Then apply the appropriate techniques to construct a rigorous argument. The solution must be self-contained and reference only the material from the course.

\vspace{0.5cm}

\textbf{Solution 7:}

The edge case analysis reveals the subtle assumptions underlying standard treatments of Shakespeare. By constructing explicit counterexamples, we see that the usual hypotheses cannot be weakened without loss of the theorem's conclusion.

\vspace{0.5cm}

\textbf{Solution 8:}

Solution approach: First recall the definitions and key theorems related to Modern Drama from 21L.005. Then apply the appropriate techniques to construct a rigorous argument. The solution must be self-contained and reference only the material from the course.

\vspace{0.5cm}

\textbf{Solution 9:}

The proof follows from the definitions and properties established in 21L.005. The key insight is to apply the relevant theorem and verify the hypotheses hold. Detailed solution depends on the specific formulation of the theorem.

\vspace{0.5cm}

\textbf{Solution 10:}

Solution approach: First recall the definitions and key theorems related to Theatricality from 21L.005. Then apply the appropriate techniques to construct a rigorous argument. The solution must be self-contained and reference only the material from the course.

\vspace{0.5cm}

\textbf{Solution 11:}

Solution approach: First recall the definitions and key theorems related to Greek Tragedy from 21L.005. Then apply the appropriate techniques to construct a rigorous argument. The solution must be self-contained and reference only the material from the course.

\vspace{0.5cm}

\textbf{Solution 12:}

The proof follows from the definitions and properties established in 21L.005. The key insight is to apply the relevant theorem and verify the hypotheses hold. Detailed solution depends on the specific formulation of the theorem.

\vspace{0.5cm}

\textbf{Solution 13:}

The proof follows from the definitions and properties established in 21L.005. The key insight is to apply the relevant theorem and verify the hypotheses hold. Detailed solution depends on the specific formulation of the theorem.

\vspace{0.5cm}

\textbf{Solution 14:}

The edge case analysis reveals the subtle assumptions underlying standard treatments of Performance Theory. By constructing explicit counterexamples, we see that the usual hypotheses cannot be weakened without loss of the theorem's conclusion.

\vspace{0.5cm}

\textbf{Solution 15:}

The proof follows from the definitions and properties established in 21L.005. The key insight is to apply the relevant theorem and verify the hypotheses hold. Detailed solution depends on the specific formulation of the theorem.

\vspace{0.5cm}

\textbf{Solution 16:}

Solution approach: First recall the definitions and key theorems related to Greek Tragedy from 21L.005. Then apply the appropriate techniques to construct a rigorous argument. The solution must be self-contained and reference only the material from the course.

\vspace{0.5cm}

\textbf{Solution 17:}

Solution approach: First recall the definitions and key theorems related to Shakespeare from 21L.005. Then apply the appropriate techniques to construct a rigorous argument. The solution must be self-contained and reference only the material from the course.

\vspace{0.5cm}

\textbf{Solution 18:}

The proof follows from the definitions and properties established in 21L.005. The key insight is to apply the relevant theorem and verify the hypotheses hold. Detailed solution depends on the specific formulation of the theorem.

\vspace{0.5cm}

\textbf{Solution 19:}

The proof follows from the definitions and properties established in 21L.005. The key insight is to apply the relevant theorem and verify the hypotheses hold. Detailed solution depends on the specific formulation of the theorem.

\vspace{0.5cm}

\textbf{Solution 20:}

The proof follows from the definitions and properties established in 21L.005. The key insight is to apply the relevant theorem and verify the hypotheses hold. Detailed solution depends on the specific formulation of the theorem.

\vspace{0.5cm}

\textbf{Solution 21:}

Solution approach: First recall the definitions and key theorems related to Greek Tragedy from 21L.005. Then apply the appropriate techniques to construct a rigorous argument. The solution must be self-contained and reference only the material from the course.

\vspace{0.5cm}

\textbf{Solution 22:}

The edge case analysis reveals the subtle assumptions underlying standard treatments of Shakespeare. By constructing explicit counterexamples, we see that the usual hypotheses cannot be weakened without loss of the theorem's conclusion.

\vspace{0.5cm}

\textbf{Solution 23:}

The proof follows from the definitions and properties established in 21L.005. The key insight is to apply the relevant theorem and verify the hypotheses hold. Detailed solution depends on the specific formulation of the theorem.

\vspace{0.5cm}

\textbf{Solution 24:}

The proof follows from the definitions and properties established in 21L.005. The key insight is to apply the relevant theorem and verify the hypotheses hold. Detailed solution depends on the specific formulation of the theorem.

\vspace{0.5cm}

\textbf{Solution 25:}

The edge case analysis reveals the subtle assumptions underlying standard treatments of Theatricality. By constructing explicit counterexamples, we see that the usual hypotheses cannot be weakened without loss of the theorem's conclusion.

\vspace{0.5cm}

\end{document}